\section{Reference metric and regular remainder}
\label{sec:ledger}

The final reference metric $\Gfr_T$ is the global flat-lag metric defined in
\eqref{eq:global-reference-Gram}. For the regular-remainder estimates it is
convenient to use a shell-local auxiliary metric. Let
$\mathcal W_T^{\rm loc}$ be the composition of
\begin{enumerate}[label=(\roman*)]
\item packet synthesis followed by the logarithmic unitary and the fixed
bounded-overlap shell windows, with the input weight $x^{-1/2}$ appearing in
\eqref{eq:local-primitive-input};
\item the invariant-source specialization of
Proposition~\ref{prop:trivial-outer-twist}; the outer translation and divisor
maps are the identity and introduce no additional arithmetic fibre;
\item the local primitive $q_j\mapsto Q_j$ of
Lemma~\ref{lem:global-local-triangular}, with the endpoint mean constraint imposed
only after restriction to $V_{\rm mean}$;
\item multiplication of the primitive output by
$\sqrt2\,L^{-1}x^{-1/2}$, so that its ordinary $L^2(dx)$ norm is exactly the
weighted norm in \eqref{eq:local-triangular-shell}, together with the remaining
positive weights common to both normalizations.
\end{enumerate}
Put
\begin{equation}
\boxed{
\Gfr_T^{\rm loc}:=(\mathcal W_T^{\rm loc})^*\mathcal W_T^{\rm loc}.
}
\label{eq:local-reference-Gram}
\end{equation}
The stationary Fourier restriction is not part of this reference map; it is an
orthogonal compression applied to the source operators and can only decrease
their Hilbert--Schmidt norms. The signed coefficient $G(\rho)$ is part of neither reference metric. By
Proposition~\ref{prop:trivial-outer-twist} there is no nontrivial outer
curvature deformation $d^{p/L}$ in the invariant source.

Lemma~\ref{lem:global-local-triangular}, applied to the common packet/core map,
gives on $V_{\rm mean}$
\begin{equation}
\boxed{
\Gfr_T^{\rm loc}\preceq C_\triangle\Gfr_T.
}
\label{eq:local-global-Gram-comparison}
\end{equation}
Both forms are positive definite on every finite packet-coordinate compression.
Indeed, if $\mathcal W_T^{\rm loc}v=0$, then every local primitive $Q_j$
vanishes; differentiation gives $q_j=0$. Since $\sum_j\psi_j^2=1$, the
logarithmic-unitary output is zero on the open physical interval. Hence the
underlying packet synthesis
\[
u\longmapsto \sum_\nu v_\nu e^{i\tau_\nu u}
\]
vanishes on a nonempty open interval. It is a finite exponential polynomial;
analytic uniqueness (equivalently the Vandermonde argument of
Lemma~\ref{lem:evaluation-rank}) gives $v=0$. The global map is injective by the
same argument: vanishing of its primitive makes the packet synthesis zero on
the global interval. The mean-zero conditions are needed for the endpoint
identities, not for injectivity.

\begin{lemma}[Polynomial floor for the core reference metric]
\label{lem:core-reference-floor}
On $V_{\rm core}\cap\ker\mathcal M_T^{\rm glob}$ there is a fixed exponent
$C_G$ such that
\begin{equation}
\boxed{
\Gfr_T\succeq T^{-2}L^{-C_G}I.
}
\label{eq:core-reference-floor}
\end{equation}
\end{lemma}

\begin{proof}
By \eqref{eq:explicit-flat-analysis}, after translating the flat core to
$0\le u\le W_T$ the global reference input is exactly
\[
g(u)=L_{\rm pkt}^{-1/2}\sum_\nu
\widetilde v_\nu e^{i\tau_\nu u},
\qquad
\widetilde v_\nu=e^{i\tau_\nu u_0}v_\nu.
\]
The phase $v\mapsto\widetilde v$ is unitary and commutes with
$D_\tau v=(\tau_\nu^{-1}v_\nu)_\nu$. Hence the two inequalities in
Lemma~\ref{lem:stable-core-frame}, after translating the integration interval,
apply directly to this synthesis and to its $D_\tau$ multiple.

Write
\[
R(u)=\int_0^u g(s)\,ds
=a+L_{\rm pkt}^{-1/2}\sum_\nu
\frac{\widetilde v_\nu}{i\tau_\nu}e^{i\tau_\nu u}.
\]
Since $T\le\tau_\nu\le2T+O(1)$, Lemma~\ref{lem:stable-core-frame} gives
\[
\left\|L_{\rm pkt}^{-1/2}\sum_\nu
\frac{\widetilde v_\nu}{i\tau_\nu}e^{i\tau_\nu u}
\right\|_{L^2(0,W_T)}^2
\gg T^{-2}\|v\|_2^2.
\]
The constant $a$ cannot cancel this oscillatory norm. Indeed
$|\int_0^{W_T}e^{i\tau_\nu u}du|\le2/|\tau_\nu|$; using
$d_T=O(TL)$ and $L_{\rm pkt}\asymp W_T\asymp L$, the squared norm of the
projection of the oscillatory sum onto the constants is
$O(T^{-3})\|v\|_2^2$. Hence
\[
\|R\|_{L^2(0,W_T)}^2\gg T^{-2}\|v\|_2^2.
\]
Inserting the factor $2L^{-2}$ in the normalized triangular primitive gives
\eqref{eq:core-reference-floor}, after enlarging the fixed exponent $C_G$.
\end{proof}

The final perturbation will use only fixed-$T$ compactness, but
\eqref{eq:core-reference-floor} is used for the direct terminal-tail estimate.
The reference metric itself contains no curvature parameter.  The finite-$T$
$p$-dependent deformation of the completed model belongs instead to
$R_T^{\rm par}$ and the AFE/stationary regular classes estimated below.

The factor used in the shell normalization is fixed by the change from the
normalized flat lag coordinate to the physical logarithmic coordinate. If
$r=u/L$ and \(\widehat q(r)=L^{1/2}g(Lr)\), then for every lag profile $H$,
\begin{equation}
\iint H(1-|r-s|)\widehat q(r)\overline{\widehat q(s)}\,dr\,ds
=\frac1L\iint
H\!\left(1-\frac{|u-v|}{L}\right)
g(u)\overline{g(v)}\,du\,dv.
\label{eq:flat-lag-jacobian}
\end{equation}
After $x=e^u$, the logarithmic unitary and the input factor
$x^{-1/2}$ in \eqref{eq:local-primitive-input} cancel the measure Jacobians.
Thus every normalized local lag term carries the same external factor $L^{-1}$;
all remaining height, stationary, and fixed local-source factors are contained
in its smooth operator-valued amplitude.

\begin{lemma}[Shell transport of the triangular normalization]
\label{lem:shell-normalization-transfer}
Let $I_j=[J_j,J_j+1]$, $x=J_j+t$, $y=J_j+s$, and let
\[
\rho_j(t,s)=\nu-\frac{|\log x-\log y|}{L},
\qquad |\nu-1|\ll L^{-1}.
\]
Suppose a Hilbert-valued shell kernel is a finite sum of terms
\begin{equation}
D_j(t,s)=\frac1L\,a_j(t,s)H(\rho_j(t,s)),
\label{eq:normalized-shell-kernel}
\end{equation}
where, for integers $a,b\ge0$ with $a+b\le2$,
\[
\|\partial_t^a\partial_s^b a_j(t,s)\|
\ll L^C J_j^{-a-b},
\]
and $H,H',H''$ are bounded by a fixed power of $L$ on the relevant compact
interval.  Here $H$ may be scalar-valued or may take values in the bounded
operators on an auxiliary Hilbert output fibre, with the displayed derivative
bounds interpreted in operator norm. Then the diagonal derivative jump satisfies
\begin{equation}
\boxed{
J_{D_j}(t)=\omega_j(t)\,a_j(t,t)H'(\nu),
\qquad
\omega_j(t)=\frac{2}{L^2(J_j+t)},
}
\label{eq:shell-jump-normalization}
\end{equation}
with the convention for $J_D$ in Lemma~\ref{lem:two-variable-primitive}.
Moreover the regular
mixed derivative obeys
\begin{equation}
\boxed{
\left\|\omega_j(t)^{-1/2}
D_{j,ts}^{\rm reg}(t,s)
\omega_j(s)^{-1/2}\right\|
\ll L^{C'}
}
\label{eq:shell-mixed-normalization}
\end{equation}
for a fixed $C'$. The same conclusions hold for finite sums with finite-dimensional
operator-valued amplitudes $a_j$.
\end{lemma}

\begin{proof}
Off the diagonal,
\[
\partial_t\rho_j
=-\frac{\operatorname{sgn}(t-s)}{Lx},
\qquad
\partial_s\rho_j
=\frac{\operatorname{sgn}(t-s)}{Ly}.
\]
The derivatives of the smooth amplitude $a_j$ have no diagonal jump. The two
one-sided values of $\partial_tH(\rho_j)$ therefore differ by
$2H'(\nu)/(Lx)$, and the external factor $L^{-1}$ in
\eqref{eq:normalized-shell-kernel} gives
\eqref{eq:shell-jump-normalization}.

For the regular mixed derivative, Leibniz' rule produces terms with two
amplitude derivatives, one amplitude and one lag derivative, or two lag
derivatives. On a length-one shell $x,y\asymp J_j$; hence they are bounded,
respectively, by fixed polylogarithmic multiples of
\[
L^{-1}J_j^{-2},
\qquad L^{-2}J_j^{-2},
\qquad L^{-3}J_j^{-2}.
\]
Since
\[
\omega_j(t)^{1/2}\omega_j(s)^{1/2}\asymp L^{-2}J_j^{-1},
\]
division by the reference weights leaves at most a fixed factor
$L^{C'}$ (indeed the first class costs only $O(L/J_j)$). This proves
\eqref{eq:shell-mixed-normalization}. The calculation is purely algebraic in the
output operators, so it is unchanged for operator-valued $H$ and amplitudes;
after finite compression one may take operator or Hilbert--Schmidt norms.
\end{proof}

\begin{lemma}[Two-variable primitive remainder criterion]
\label{lem:two-variable-primitive}
Let $I=[0,1]$ and let $q=Q'$ with $Q(0)=Q(1)=0$. Let $D(x,y)$ be a Hermitian
kernel which is continuous across the diagonal, $C^1$ on each open triangle,
and whose regular mixed derivative is in $L^2(I^2)$. Define the diagonal jump
\[
J_D(x):=\partial_xD(x,x+)-\partial_xD(x,x-).
\]
Then, distributionally,
\[
\partial_x\partial_yD
=J_D(x)\delta_{x=y}+D_{xy}^{\rm reg},
\]
and two integrations by parts give
\begin{equation}
\boxed{
\iint D(x,y)q(x)\overline{q(y)}\,dxdy
=\int J_D(x)|Q(x)|^2dx
+\iint D_{xy}^{\rm reg}(x,y)Q(x)\overline{Q(y)}\,dxdy.
}
\label{eq:two-variable-primitive}
\end{equation}
For the global application, let $D_{j,T}$ be the corresponding
operator-valued kernel on the finite output fibre of the $j$th retained unit
shell. All common positive output weights are already included in
$\mathcal W_T^{\rm loc}$; by Proposition~\ref{prop:trivial-outer-twist} there is
no additional outer arithmetic metric. In the local coordinate $0\le t\le1$
put
\begin{equation}
\omega_j(t):=\frac{2}{L^2(J_j+t)}.
\label{eq:local-primitive-weight}
\end{equation}
By \eqref{eq:local-triangular-shell}, $\omega_j(t)$ is the exact scalar
primitive reference density on this output fibre. Define the normalized jump
and regular mixed kernel by
\[
\mathcal J_{j,T}(t)
:=\omega_j(t)^{-1}J_{D_{j,T}}(t),
\]
\[
\mathcal K_{j,T}(t,s)
:=\omega_j(t)^{-1/2}D_{j,T,ts}^{\rm reg}(t,s)
\omega_j(s)^{-1/2}.
\]
If
\begin{equation}
\sup_{j,x}\|\mathcal J_{j,T}(x)\|_{\op}
=\delta_T=o(1)
\label{eq:jump-criterion}
\end{equation}
and
\begin{equation}
M_T:=\sum_{j\in\mathcal J_T}\iint
\|\mathcal K_{j,T}(x,y)\|_{S_2}^2\,dxdy=o(N_T),
\label{eq:mixed-criterion}
\end{equation}
then every finite retained compression of total dimension $O(N_T)$ satisfies
\begin{equation}
\boxed{
\|\Gfr_T^{-1/2}R_T\Gfr_T^{-1/2}\|_{S_2}^2
\ll \delta_T^2N_T+M_T=o(N_T).
}
\label{eq:primitive-HS-criterion}
\end{equation}
In particular, a uniform per-shell estimate
\[
\iint\|\mathcal K_{j,T}(x,y)\|_{S_2}^2\,dxdy\ll L^C
\]
for fixed $C$ implies \eqref{eq:mixed-criterion} whenever
$A_0>C+2$, since $\#\mathcal J_T=O(X_T)$ and
$X_T=T/(2\pi L^{A_0})$.
\end{lemma}

\begin{proof}
The endpoint conditions kill all boundary terms in the two integrations by
parts, proving \eqref{eq:two-variable-primitive}. On the $j$th local primitive
space the exact reference quadratic form is
\[
\int_0^1\omega_j(t)\|Q(t)\|^2\,dt.
\]
Conjugating by multiplication with $\omega_j(t)^{1/2}$ turns the diagonal term
in \eqref{eq:two-variable-primitive} into the multiplication operator
$M_{\mathcal J_{j,T}}$ and the regular term into the integral operator
$T_{\mathcal K_{j,T}}$. On the orthogonal direct sum of shell outputs put
\[
B_{\rm jump}:=\bigoplus_jM_{\mathcal J_{j,T}},
\qquad
B_{\rm reg}:=\bigoplus_jT_{\mathcal K_{j,T}}.
\]
Then
\[
\|B_{\rm jump}\|_{\op}\le\delta_T,
\qquad
\|B_{\rm reg}\|_{S_2}^2
\le\sum_j\iint\|\mathcal K_{j,T}(t,s)\|_{S_2}^2\,dt\,ds=M_T.
\]
The multiplication operator $B_{\rm jump}$ need not itself be
Hilbert--Schmidt on the continuous shell-output space; finite dimensionality
enters only after pullback. Put $G_{\rm loc}:=\Gfr_T^{\rm loc}$ on the retained
compression, $W:=\mathcal W_T^{\rm loc}$, and
\[
V:=WG_{\rm loc}^{-1/2}.
\]
Then $V$ is an isometry from the retained packet space, of dimension
$d_T=O(N_T)$, into $\Ran W$. Therefore
\[
\|V^*B_{\rm jump}V\|_{S_2}^2
\le d_T\|B_{\rm jump}\|_{\op}^2
\ll\delta_T^2N_T,
\]
whereas the ideal property of the Hilbert--Schmidt norm gives
\[
\|V^*B_{\rm reg}V\|_{S_2}
\le\|B_{\rm reg}\|_{S_2}\le M_T^{1/2}.
\]
Since all common output weights are already contained in
$\mathcal W_T^{\rm loc}$, the pulled-back remainder is
$V^*(B_{\rm jump}+B_{\rm reg})V$, and we obtain
\[
\|G_{\rm loc}^{-1/2}R_TG_{\rm loc}^{-1/2}\|_{S_2}^2
\ll\delta_T^2N_T+M_T.
\]
Finally \eqref{eq:local-global-Gram-comparison} gives
$G_{\rm loc}\preceq C_\triangle\Gfr_T$. With
\[
A:=G_{\rm loc}^{1/2}\Gfr_T^{-1/2}
\]
one has $\|A\|_{\op}^2\le C_\triangle$ and
\[
\Gfr_T^{-1/2}R_T\Gfr_T^{-1/2}
=A^*\bigl(G_{\rm loc}^{-1/2}R_TG_{\rm loc}^{-1/2}\bigr)A.
\]
Thus global whitening costs only the fixed factor $C_\triangle$, proving
\eqref{eq:primitive-HS-criterion}.
\end{proof}

\begin{lemma}[Reference compatibility of the common HLZ output]
\label{lem:HLZ-output-reference-compatibility}
Restrict to the stable core and to $V_{\rm mean}$ of
\eqref{eq:mean-map}.  Let $\mathscr C_T$ be the shift-independent common
one-sided output analysis of Lemma~\ref{lem:model-HLZ-Fubini-lift}, with the
same normalized one-sided factors as in the definition of $\Gfr_T$.  After
forming the nonnegative height-block direct sum with weights
$\omega_b^{1/2}$, put
\[
G_{\rm out}:=\mathscr C_T^*\mathscr C_T.
\]
On every finite retained packet compression there is a fixed exponent
$C_{\rm out}$, independent of the compression dimension, such that
\begin{equation}
\boxed{
G_{\rm out}\preceq L^{C_{\rm out}}\Gfr_T.
}
\label{eq:HLZ-output-reference-comparison}
\end{equation}
The same statement holds after adjoining any of the fixed packet-independent
finite source, curvature, or Cauchy output fibres occurring in the finite model.
\end{lemma}

\begin{proof}
For one height block and one packet pair, the unit-multiplier output Gram of
Lemma~\ref{lem:model-HLZ-Fubini-lift} is obtained from the intersection of the
two translated one-sided prefix intervals.  In the normalized common-prefix
coordinate its scalar overlap is therefore
\[
\rho_{rs}^{(T)}
=\nu_T-\frac{|w_r-w_s|}{L}
=\rho_{rs}^{(0)}+\delta(\tau),
\]
by \eqref{eq:finiteT-common-endpoint}.  The second term is independent of the
packet pair.  Consequently its quadratic form factors through the global mean
map $\mathcal M_T^{\rm glob}$ and vanishes on $V_{\rm mean}$.

The remaining kernel is
\[
\rho_{rs}^{(0)}=1-\frac{|w_r-w_s|}{L},
\]
namely the flat triangular lag kernel.  Under the exact principal dilation
\eqref{eq:dilation-covariance} and the normalized common-core identification
\eqref{eq:normalized-common-core-analysis}, this is precisely the feature
factorization of the triangular primitive form $Q_\triangle$ used to define
$\Gfr_T$ in \eqref{eq:global-reference-Gram}; on the mean-zero space the
identity
\[
\iint(1-|r-s|)q(r)\overline{q(s)}\,dr\,ds
=2\int\left|\int_0^rq(v)\,dv\right|^2dr
\]
is exact.  Thus the unit common-prefix feature introduces no packet-dimension
loss.

The remaining normalized stationary, archimedean, block, and finite Cauchy
factors are packet-label independent and, through the fixed derivative orders
used here, cost at most a fixed power of $L$.  Since the height partition is
nonnegative and $\sum_b\omega_b=1$, using $\omega_b^{1/2}$ in the orthogonal
block direct sum reproduces the linear block sum without a cross term.
This proves \eqref{eq:HLZ-output-reference-comparison}.  Tensoring with a
fixed packet-independent finite output fibre changes only the fixed exponent,
not the packet dimension.
\end{proof}

\begin{lemma}[Operator lift of the finite-model HLZ contour error]
\label{lem:HLZ-diagonal-operator-lift}
Choose the arbitrary logarithmic exponent $A=A_{\rm diag}$ in
Lemma~\ref{lem:packet-HLZ-diagonal} large enough to absorb all fixed packet,
shift, curvature, shell, and Cauchy-residue derivative losses, and put
\[
\varepsilon_{\rm diag}
:=L^{-A_{\rm diag}+C_{\rm diag}+o(1)}=o(1).
\]
Let $R_T^{\rm com}$ denote the scalar HLZ diagonalization error of the finite
model source after the physical-shift Cauchy residues. After restriction to the
stable core and compression to any retained packet space contained in
$V_{\rm mean}$ and of dimension $d_T=O(N_T)$,
\begin{equation}
\boxed{
\left\|\Gfr_T^{-1/2}R_T^{\rm com}\Gfr_T^{-1/2}\right\|_{S_2}^2
\ll N_T\varepsilon_{\rm diag}^2=o(N_T).
}
\label{eq:HLZ-diagonal-relative-HS}
\end{equation}
No packet-dimension factor is lost before the final Hilbert--Schmidt conversion.
\end{lemma}

\begin{proof}
Lemma~\ref{lem:model-HLZ-Fubini-lift} gives the error before packet pullback in
the form of a packet-label-independent multiplication operator on the same
one-sided output fibre as the finite model main term.  Its multiplier
$\varepsilon_T$ satisfies
\[
\|\varepsilon_T\|_{\infty}
\ll\varepsilon_{\rm diag}
\]
after all fixed derivatives and finite Cauchy residues.  Let
$\mathscr C_T$ be the normalized output analysis from
Lemma~\ref{lem:model-HLZ-Fubini-lift} and put
$G_{\rm out}=\mathscr C_T^*\mathscr C_T$.  Then on the support of
$G_{\rm out}$
\[
\|G_{\rm out}^{-1/2}\mathscr C_T^*M_{\varepsilon_T}\mathscr C_T
G_{\rm out}^{-1/2}\|_{\op}
\le\|\varepsilon_T\|_\infty.
\]
By Lemma~\ref{lem:HLZ-output-reference-compatibility},
$G_{\rm out}\preceq L^{C_{\rm out}}\Gfr_T$ on the retained mean-zero space.
Using the support inverse of $G_{\rm out}$, put
\[
A:=G_{\rm out}^{1/2}\Gfr_T^{-1/2}.
\]
Then $\|A\|_{\op}^2\ll L^{C_{\rm out}}$ and, because
$R_T^{\rm com}=\mathscr C_T^*M_{\varepsilon_T}\mathscr C_T$ has range in the
support of $G_{\rm out}$,
\[
\Gfr_T^{-1/2}R_T^{\rm com}\Gfr_T^{-1/2}
=A^*\bigl(G_{\rm out}^{-1/2}R_T^{\rm com}G_{\rm out}^{-1/2}\bigr)A.
\]
Hence its operator norm is
$O(L^{C_{\rm out}}\varepsilon_{\rm diag})$.  Enlarge
$C_{\rm diag}$ to absorb this fixed output-reference loss and choose
$A_{\rm diag}$ accordingly. Compression to $d_T=O(N_T)$ dimensions and
$\|B\|_{S_2}^2\le d_T\|B\|_{\op}^2$ then give
\eqref{eq:HLZ-diagonal-relative-HS}.  Thus the packet dimension enters only in
the final Hilbert--Schmidt conversion, not in the scalar-to-operator lift.  No
levelwise exact-HLP Perron deformation is used in this route.
\end{proof}

\begin{lemma}[Finite-$T$ parameter deformation is relative-regular]
\label{lem:finiteT-parameter-regular}
Let $R_T^{\rm par}$ be the finite-$T$ parameter defect in
\eqref{eq:common-core-principal-congruence}.  After compression to the retained
mean-zero space, its normalized diagonal jump and regular mixed kernel satisfy
\begin{equation}
\boxed{
\sup_{j,t}\|\mathcal J^{\rm par}_{j,T}(t)\|_{\op}\ll L^{-1}=o(1),
\qquad
\iint\|\mathcal K^{\rm par}_{j,T}(t,s)\|_{S_2}^2\,dt\,ds
\ll L^{C_{\rm par}}
}
\label{eq:finiteT-parameter-shell-bounds}
\end{equation}
for one fixed exponent $C_{\rm par}$.  Consequently
\begin{equation}
\boxed{
\left\|\Gfr_T^{-1/2}R_T^{\rm par}\Gfr_T^{-1/2}\right\|_{S_2}^2
=o(N_T).
}
\label{eq:finiteT-parameter-relative-HS}
\end{equation}
No blockwise mean-zero condition and no scalar variable-coefficient kernel in
the final lag variables is used.
\end{lemma}

\begin{proof}
Put
\[
H_{b,\tau}^{\rm par}(\rho)
:=[p^2]\,\Delta_{b,\tau}(\rho;p),
\]
with $\Delta_{b,\tau}$ from \eqref{eq:finiteT-parameter-defect}.  The derivative
bound \eqref{eq:finiteT-parameter-defect-derivatives} gives
\begin{equation}
\max_{0\le a\le2}\sup_{b,\tau,\rho}
|\partial_\rho^aH_{b,\tau}^{\rm par}(\rho)|\ll L^{-1}.
\label{eq:finiteT-parameter-profile-small}
\end{equation}
We now work shell by shell and keep the physical height inside the ordinary
smooth stationary amplitude; we do not promote it to a new output coordinate.
For one retained shell and one of the finitely many local source/refinement
labels, let $\tau_{j,\alpha}(t,s)\asymp T$ denote the physical height sampled by
the principal stationary parametrization.  Before differentiating, sum the
slow height blocks at that physical height and set
\begin{equation}
\mathcal H^{\rm par}_{j,\alpha,T}(t,s;\rho)
:=\sum_b\omega_b\!\left(\tau_{j,\alpha}(t,s)\right)
 H^{\rm par}_{b,\tau_{j,\alpha}(t,s)}(\rho).
\label{eq:finiteT-parameter-shell-profile}
\end{equation}
The partition was inserted linearly in the contour integral, so this is the
actual block recombination: there are no products
$\omega_b\omega_{b'}$.  Since the weights are nonnegative and sum to one,
\eqref{eq:finiteT-parameter-profile-small} gives, uniformly in the shell
variables and labels,
\begin{equation}
\boxed{
\max_{0\le a\le2}
\|\partial_\rho^a\mathcal H^{\rm par}_{j,\alpha,T}(t,s;\rho)\|_{\op}
\ll L^{-1}.
}
\label{eq:finiteT-parameter-operator-profile}
\end{equation}
This statement is local in the stationary variables and makes no assertion
that physical height is a function only of the final lag.

After the logarithmic unitary and common stationary normalization, collect the
fixed local labels in the existing finite shell output fibre.  The parameter
defect is a finite sum of terms
\begin{equation}
D_{j,T}^{\rm par}(t,s)
=\frac1L\,a_{j,T}(t,s)
 \mathcal H^{\rm par}_{j,T}\!\left(t,s;
 1-\frac{|\log(J_j+t)-\log(J_j+s)|}{L}\right),
\label{eq:finiteT-parameter-shell-kernel}
\end{equation}
with all common positive one-sided factors already included in the local
reference map.  The dependence of $\tau_{j,\alpha}$ and of the block weights on
$t,s$ is smooth.  Lemma~\ref{lem:stationary-symbol},
\eqref{eq:principal-field-symbol-bound}, and
\eqref{eq:height-partition-derivatives}, followed by the chain rule, give
through the two derivatives used by the primitive identity
\begin{equation}
\|\partial_t^a\partial_s^b a_{j,T}(t,s)\|_{\op}
+\|\partial_t^a\partial_s^b
  \mathcal H^{\rm par}_{j,T}(t,s;\rho)\|_{\op}
\ll L^{C_{\rm par}}J_j^{-a-b},
\qquad a+b\le2,
\label{eq:finiteT-parameter-amplitude-symbol}
\end{equation}
uniformly for $\rho$ in the relevant compact interval.  Only a fixed
polylogarithmic bound is asserted for the $t,s$ derivatives; their smallness is
not needed.

There is one stronger zeroth-order fact on the diagonal.  After the common
positive one-sided weights are placed in $\mathcal W_T^{\rm loc}$ as in
\eqref{eq:local-reference-Gram}, the remaining relative amplitude satisfies
\begin{equation}
\boxed{
\sup_{j,t}\|a_{j,T}(t,t)\|_{\op}\ll1.
}
\label{eq:finiteT-parameter-diagonal-amplitude}
\end{equation}
Indeed $R_T^{\rm par}$ changes only the scalar completed profile
\eqref{eq:finiteT-parameter-defect}; its two one-sided packet/stationary factors
are the same as those used in the reference output.

All $t,s$ dependence in $a_{j,T}$,
$\tau_{j,\alpha}$, and $\omega_b(\tau_{j,\alpha})$ is smooth across the
diagonal.  Hence their first derivatives have the same two one-sided limits.
The only source of a diagonal derivative jump in
\eqref{eq:finiteT-parameter-shell-kernel} is therefore the derivative of the
absolute lag.  The calculation in
Lemma~\ref{lem:shell-normalization-transfer} gives, after division by the exact
reference density,
\[
\|\mathcal J^{\rm par}_{j,T}(t)\|_{\op}
\ll
\|a_{j,T}(t,t)\|_{\op}
\|\partial_\rho\mathcal H^{\rm par}_{j,T}(t,t;1)\|_{\op}
\ll L^{-1}.
\]
This is the first estimate in \eqref{eq:finiteT-parameter-shell-bounds}.  For
the regular mixed derivative, Leibniz' rule together with
\eqref{eq:finiteT-parameter-operator-profile} and
\eqref{eq:finiteT-parameter-amplitude-symbol} gives the second estimate in
\eqref{eq:finiteT-parameter-shell-bounds}.  In this step only a fixed
polylogarithmic per-shell bound is required, not an $O(L^{-1})$ bound.

Apply Lemma~\ref{lem:two-variable-primitive}.  The jump part contributes at
most
\[
O(L^{-2}N_T)=o(N_T)
\]
to the squared relative Hilbert--Schmidt norm.  The regular mixed part is
bounded by
\[
\sum_{j\in\mathcal J_T}L^{C_{\rm par}}
\ll X_TL^{C_{\rm par}}
=TL^{C_{\rm par}-A_0+O(1)}=o(N_T),
\]
after enlarging the fixed regular exponent $C_{\rm reg}$, before $A_0$ is
chosen, so that it dominates $C_{\rm par}$ in
\eqref{eq:acyclic-A0-choice}.  Finally
\eqref{eq:local-global-Gram-comparison} transfers the shell-local estimate to
$\Gfr_T$ with only the fixed factor $C_\triangle$.  This proves
\eqref{eq:finiteT-parameter-relative-HS}.
\end{proof}

\begin{lemma}[The critical-$J$ AFE and polar discrepancy is relative-regular]
\label{lem:J-AFE-relative-regular}
Let $R_{J,\rm AFE}$ denote the difference between the exact smoothed
critical-$J$ self contribution and the sharp full-fibre contribution retained
in the frozen model.  Here ``exact'' includes all four remote residues
$u=-s_1,1-s_1,-s_2,1-s_2$ in
\eqref{eq:AFE-four-polar-residues}, with their sign in
\eqref{eq:exact-shifted-product-AFE}.  Put
\[
\varepsilon_{\rm AFE}:=L^3/T+T^{-10}.
\]
On the retained mean-zero space,
\begin{equation}
\boxed{
\sup_{j,t}\|\mathcal J^{J,\rm AFE}_{j,T}(t)\|_{\op}
\ll\varepsilon_{\rm AFE}=o(1),
\qquad
\iint\|\mathcal K^{J,\rm AFE}_{j,T}(t,s)\|_{S_2}^2\,dt\,ds
\ll L^{C_{\rm AFE}}
}
\label{eq:J-AFE-shell-regular}
\end{equation}
for a fixed exponent $C_{\rm AFE}$.  Consequently
\begin{equation}
\boxed{
\left\|\Gfr_T^{-1/2}R_{J,\rm AFE}\Gfr_T^{-1/2}\right\|_{S_2}^2
=o(N_T).
}
\label{eq:J-AFE-relative-HS}
\end{equation}
\end{lemma}

\begin{proof}
Work on one retained shell after the common logarithmic unitary, but keep every
physical-height factor inside the ordinary smooth stationary amplitude.  For a
fixed local source/refinement label let
$\tau_{j,\alpha}(t,s)\asymp T$ be the physical height sampled by the principal
stationary parametrization, and define the local AFE error profile by inserting
that height into the two-variable remainder
$\mathcal R^J_{U,T}(A,B;\tau)$ of
\eqref{eq:J-AFE-sharp-transfer} before the shell lag is substituted.  This
remainder includes the four Gaussian-small polar terms.  After normalized
curvature extraction write the resulting profile as
\[
\mathcal H^{J,\rm AFE}_{j,\alpha,T}(t,s;\rho;A,B,\tau).
\]
Lemma~\ref{lem:J-AFE-sharp-transfer}, in particular the joint gamma-ratio,
dual-multiplier, and polar estimates
\eqref{eq:AFE-joint-gamma-ratio}, \eqref{eq:AFE-joint-dual-factor}, and
\eqref{eq:AFE-polar-Gaussian-bound}, gives uniformly in the shell variables
\begin{equation}
\boxed{
\max_{a+b+c+m\le2}
\|\partial_A^a\partial_B^b\partial_\rho^c(T\partial_\tau)^m
 \mathcal H^{J,\rm AFE}_{j,\alpha,T}
 (t,s;\rho;A,B,\tau)\|_{\op}
\ll\varepsilon_{\rm AFE}.
}
\label{eq:J-AFE-local-profile-small}
\end{equation}
Here the $\rho$ derivatives are the independent endpoint derivatives in
\eqref{eq:J-AFE-sharp-error}.  On imposing
$\tau=2\pi e^{L\rho}$, total endpoint differentiation is
$\partial_\rho+L\tau\partial_\tau$; through second order the sharper
$O(L^j/T)$ estimate in \eqref{eq:J-AFE-sharp-error} is still bounded by
$\varepsilon_{\rm AFE}$.  No characteristic-function replacement is made
inside the shell argument, and curvature derivatives are allowed to hit
$A$, $B$, the gamma ratios, and the dual multiplier separately.

Collecting the fixed local labels in the existing shell output fibre, the AFE
smoothing discrepancy is a finite sum of kernels
\begin{equation}
D^{J,\rm AFE}_{j,T}(t,s)
=\frac1L\,a^J_{j,T}(t,s)
 \mathcal H^{J,\rm AFE}_{j,T}\!\left(t,s;
 1-\frac{|\log(J_j+t)-\log(J_j+s)|}{L};
 A,B,\tau_{j,\alpha}(t,s)\right).
\label{eq:J-AFE-shell-kernel}
\end{equation}
The two one-sided packet/stationary factors are exactly those of the sharp
critical-$J$ term.  Hence, after the common positive weights have been placed in
$\mathcal W_T^{\rm loc}$, their relative diagonal amplitude is uniformly
bounded as in \eqref{eq:finiteT-parameter-diagonal-amplitude}.  Moreover the
$t,s$ dependence of the physical height, the height cutoff, and the stationary
amplitude is smooth across the diagonal.  Lemma~\ref{lem:stationary-symbol},
\eqref{eq:height-partition-derivatives}, and the chain rule give fixed
polylogarithmic bounds for their first two $t,s$ derivatives.

Consequently the only one-sided diagonal jump in
\eqref{eq:J-AFE-shell-kernel} comes from differentiating the absolute lag.
The calculation of Lemma~\ref{lem:shell-normalization-transfer}, together with
\eqref{eq:J-AFE-local-profile-small}, gives
\[
\|\mathcal J^{J,\rm AFE}_{j,T}(t)\|_{\op}
\ll\varepsilon_{\rm AFE}.
\]
For the regular mixed derivative, Leibniz' rule uses the same fixed stationary
symbol bounds and the endpoint/curvature derivatives in
\eqref{eq:J-AFE-local-profile-small}; in particular, the chain rule may
differentiate the physical height and the dual multiplier.  After reference
normalization it yields a fixed polylogarithmic per-shell Hilbert--Schmidt bound.
The four polar profiles are smaller than every power of $T$ with the same
derivatives and are therefore absorbed by the displayed $T^{-10}$ term.  This proves
\eqref{eq:J-AFE-shell-regular}.  Notice that physical height has remained a
smooth shell parameter throughout; no extra continuum output fibre and no
heightwise mean-zero constraint has been introduced.

Lemma~\ref{lem:two-variable-primitive} now bounds the squared relative norm by
\[
O(\varepsilon_{\rm AFE}^2N_T)
+O(X_TL^{C_{\rm AFE}})=o(N_T),
\]
after $C_{\rm reg}$ is chosen to dominate $C_{\rm AFE}$ before $A_0$ in
\eqref{eq:acyclic-A0-choice}.  The local-to-global Gram comparison finishes the
proof.
\end{proof}

\begin{proposition}[Relative Hilbert--Schmidt bound for the regular remainder]
\label{prop:regular-ledger}
After the finite model principal source has been extracted and the exact--model
difference has been retained as the direct residual
\eqref{eq:direct-exact-model-residual}, the regular remainder is a sum of
mutually exclusive pieces. After quadratic-form compression to $\Pret$, it
satisfies
\begin{equation}
\boxed{
\left\|
(\Gfr_T|_{\Pret})^{-1/2}(R_T^{\rm tri}|_{\Pret})
(\Gfr_T|_{\Pret})^{-1/2}
\right\|_{S_2}^2=o(N_T).
}
\label{eq:regular-HS}
\end{equation}
\end{proposition}

\begin{proof}
We treat the primitive-controlled source classes with
Lemma~\ref{lem:two-variable-primitive}; the boundary nonstationary and
high-sideband classes are estimated directly in relative Hilbert--Schmidt norm
below. Throughout, ``normalized jump'' and ``normalized mixed kernel'' mean the
quantities $\mathcal J_{j,T}$ and $\mathcal K_{j,T}$ defined there, including
the exact shell density \eqref{eq:local-primitive-weight}. The common factor
$L^{-1}$ is fixed by \eqref{eq:flat-lag-jacobian}. By
Proposition~\ref{prop:trivial-outer-twist} there is no arithmetic whitening in
the invariant source. Lemma~\ref{lem:shell-normalization-transfer} converts the
previously established lag-profile derivative bounds into these shell-normalized
quantities. All remaining fixed losses are powers of $L$ and are included in
the exponents chosen in Section~\ref{sec:retained}.

\emph{(i) No exact-HLP local-projection remainder.}
In the final source route no exact HLP level is locally projected and no
exact/leading/model telescope is inserted into the contour identity.  Hence the
former class $R_T^{\rm loc}$ is absent from the final decomposition.  The
local-Laurent comparison \eqref{eq:local-model-error} remains available only in
the optional stronger exact-HLP route and contributes nothing to
$R_T^{\rm tri}$ here.

\emph{(ii) Slow self-freezing discrepancy.}
The exact algebraic split \eqref{eq:direct-model-source-split} leaves only the
scalar self difference $f-F_b$ outside the finite model and the direct residual;
no HLP level is frozen.  By the unconditional critical-$J$ estimate
\eqref{eq:self-freezing-kernel}, the normalized jump and mixed kernel satisfy
\[
\delta_T\ll L^{-D-1+C_{\rm fr}}=o(1),
\qquad
\|\mathcal K_{\rm freeze}\|_{L^2(S_2)}
\ll L^{-D-1+C_{\rm fr}}=o(1).
\]
The bounded overlap of the smooth height partition preserves these estimates,
so \eqref{eq:primitive-HS-criterion} applies.

\emph{(iii) Hardy-gauge $f'$ source.}
This class is exactly $R_{f'}^{\rm src}$ from \eqref{eq:fprime-source}, including
the curvature, logarithmic-derivative, and quadratic $f'$ terms; none remains in
the principal source. On the high strip,
\[
f'(s)=O(T^{-1}),
\qquad |f(s)-P(s)|\asymp L
\]
on the safe right edge. The remaining factors are the same completed shifted
zeta/Cauchy factors already present in the principal master. A fixed number of
stationary, packet and curvature derivatives therefore costs only $L^C$, while
every monomial in \eqref{eq:fprime-source} retains at least one factor $T^{-1}$.
Thus
\[
\|\mathcal J_{f'}\|_{\op}\ll T^{-1}L^C=o(1),
\qquad
\|\mathcal K_{f'}\|_{L^2(S_2)}\ll T^{-1}L^C.
\]

\emph{(iv) Horizontal connectors.}
By Lemma~\ref{lem:horizontal-decay}, after the endpoint-collar deletion already
included in $V_{\rm end}$ we choose its arbitrary decay exponent as
$A=M_*+C_{\rm hor}$, where $C_{\rm hor}$ dominates the fixed polynomial losses
from the finite pullbacks and mixed derivatives. Hence
\[
\mathcal J_{\rm hor}=0,
\qquad
\|\mathcal K_{\rm hor}\|_{L^2(S_2)}\ll T^{-M_*}L^C.
\]
No further decay order is chosen later.

\emph{(v) Finite-$T$ parameter deformation, first stationary/Stirling
corrections, and the critical-$J$ AFE smooth-to-sharp and polar remainder.}
The finite-$T$ parameters $\lambda(\tau),\mu_b$ and the moving common endpoint
do not remain inside $P_T$.  Their exact sharp-model difference from the frozen
profile is $R_T^{\rm par}$ of \eqref{eq:common-core-principal-congruence}, and
Lemma~\ref{lem:finiteT-parameter-regular} gives directly
\begin{equation}
\boxed{
\left\|\Gfr_T^{-1/2}R_T^{\rm par}\Gfr_T^{-1/2}\right\|_{S_2}^2=o(N_T).
}
\label{eq:finiteT-parameter-ledger}
\end{equation}
This estimate is performed shell-locally with the physical-height dependence
retained inside the smooth stationary amplitude, so it does not use a
constant-parameter stability lemma or a global variable-coefficient lag
kernel.

The exact--model difference itself is not part of this class: it remains a
direct one-$\chi$ family and is split by sidebands in class~(vii).  The first
omitted stationary-phase and Stirling coefficients attached to the finite model
lag profiles carry an additional factor $T^{-1}L^C$ after the fixed derivatives,
by Lemma~\ref{lem:stationary-symbol} and Stirling's formula.

The other completed-model term in this class is the critical-$J$
smooth-to-sharp discrepancy together with the four remote AFE residues.
Lemma~\ref{lem:J-AFE-relative-regular}, based on
the exact four-residue AFE, explicit Mellin complement, joint gamma-ratio
estimates, and dual-multiplier derivatives of
Lemma~\ref{lem:J-AFE-sharp-transfer}, gives
\[
\left\|\Gfr_T^{-1/2}R_{J,\rm AFE}\Gfr_T^{-1/2}\right\|_{S_2}^2=o(N_T).
\]
This estimate, like the parameter-defect estimate, keeps the physical-height
dependence inside the smooth shell amplitude until the primitive normalization
has been applied.  Thus neither the finite-$T$ parameter deformation nor the
AFE smoothing/polar error is silently absorbed into the positive principal kernel.

For the stationary/Stirling part,
Lemma~\ref{lem:shell-normalization-transfer} gives
\[
\|\mathcal J_{\rm corr}\|_{\op}\ll T^{-1}L^C=o(1),
\qquad
\|\mathcal K_{\rm corr}\|_{L^2(S_2)}\ll T^{-1}L^C.
\]
The order-one complementary pole-free direct source is not included in this
class; it is treated by the sideband decomposition in (vii).

\emph{(vi) Nonstationary and smooth height-partition terms.}
The height-symbol argument in Lemma~\ref{lem:stationary-symbol}, before the
stationary expansion is taken, is uniform for the complete summed HLP
hierarchy; together with Lemma~\ref{lem:direct-carrier-square} it gives the same
fixed derivative and coefficient-square bounds for the direct residual in the
nonstationary sector.  Differentiate first. On every retained nonboundary shell
$j\in\mathcal J_T$, use the integration-by-parts order $N$ and buffer exponent
$B$ fixed in
Section~\ref{sec:retained}. On a retained boundary shell meeting
$\{1,X_T\}$, first remove the finite endpoint sum
\eqref{eq:boundary-shell-jets}; Proposition~\ref{prop:boundary-endpoint-rank}
assigns those terms only to $E_T^{\rm edge}$. The remaining interior integral
obeys the same repeated integration-by-parts estimate. By
\eqref{eq:nonstat-shell}, $R_j\ge L^B$, and the bounded-overlap shell sum,
the squared leakage before the fixed relative-normalization loss is
\[
\ll N_TL^{2C_N-B(2N-1)+o(1)}.
\]
After the primitive normalization, whose squared fixed loss is at most
$L^{C_{\rm reg}}$, the choice of $B$ and $M_*$ in
Section~\ref{sec:retained} gives
\begin{equation}
\boxed{
\mathcal J_{\rm nonstat}=0,
\qquad
\sum_j\iint
\|\mathcal K_{j,\rm nonstat}(x,y)\|_{S_2}^2\,dxdy=o(N_T).
}
\label{eq:nonstat-relative-HS}
\end{equation}
For the slow height decomposition use the partition fixed in
\eqref{eq:height-principal-partition}--\eqref{eq:height-partition-derivatives}
and insert it in the linear contour integral before stationary phase.  The
leading pieces sum exactly to the original amplitude because
$\sum_b\omega_b=1$.  Every differentiated partition factor satisfies
\[
\partial_t^r\omega_b(t)=O_r(H^{-r})
=O_r(T^{-r}L^{Dr}),
\]
so the non-leading differentiated-partition terms have zero principal cusp and
satisfy the same regular criterion.  No affine-reproduction identity is needed
in the final proof, and no block-boundary packet deletion is introduced.

\emph{(vii) High sidebands of the direct one-$\chi$ residual.}
This class contains the high sidebands of the entire direct exact--model
residual \eqref{eq:direct-exact-model-residual}, including the full
norm-summable HLP hierarchy, the ordinary $P$ term, and the translated
coefficient shift $-P(s-2/L)$, together with the other genuinely direct
stationary sources.  No exact HLP level has been locally projected.  The
auxiliary outer-excess family is zero for the invariant source.
Lemmas~\ref{lem:stationary-symbol} and \ref{lem:direct-carrier-square} cover all
the direct families occurring here. On
$I_j=[J_j,J_j+1]$ write
\[
a_{j,m}(y;p)=\sum_{r\in\mathbb Z}\widehat a_{j,m}(r;p)e(ry).
\]
Uniform fixed-order stationary phase gives, for every chosen $N$,
\[
|\partial_p^s\widehat a_{j,m}(r;p)|
\ll_{N,s}L^{C_{N,s}}(1+|r|)^{-N}.
\]
Put $d_j\asymp T/J_j$ and $R_j^{\rm sb}:=\epsilon_Ld_j$ with
$\epsilon_L=L^{-\sigma}$. The choice $A_0>\sigma+3$ gives
$R_j^{\rm sb}\ge L^3$ on every retained shell. Place
$|r|\le R_j^{\rm sb}$ in the common low-sideband edge range.

For the high sidebands choose Fourier decay order $M_{\rm sb}+4$. Let
$c_{j,m}$ be the normalized coefficient on the finite output fibre after the
fixed mixed-kernel derivatives. Since the stationary interval has upper
endpoint $B_j\asymp d_j$, Lemma~\ref{lem:direct-carrier-square} gives
\begin{equation}
\sum_{m\in\mathcal I_j}\|c_{j,m}\|_{S_2}^2
\ll d_jL^{C_{\rm sb}}.
\label{eq:carrier-square-one-shell}
\end{equation}
For the effective Fourier shift $q=m-r$ put
\[
H_{q,m}:=\widehat a_{j,m}(m-q;p)
\mathbf 1_{|q-m|>R_j^{\rm sb}}.
\]
The Fourier decay gives both row and column Schur bounds
\[
\sup_m\sum_q|H_{q,m}|+\sup_q\sum_m|H_{q,m}|
\ll L^{C_{\rm sb}}(R_j^{\rm sb})^{-M_{\rm sb}-3}.
\]
Hence the Hilbert-space-valued Schur estimate implies
\[
\sum_q\left\|\sum_mH_{q,m}c_{j,m}\right\|_{S_2}^2
\ll L^{C_{\rm sb}}d_j(R_j^{\rm sb})^{-2M_{\rm sb}-6}.
\]
Distinct compressed shifts $P_jS_qP_j$ occupy distinct matrix diagonals and are
Hilbert--Schmidt orthogonal; moreover
$\|P_jS_qP_j\|_{S_2}^2\ll d_j$. Therefore
\[
\|\mathcal K_{j,\rm hi}\|_{S_2}^2
\ll L^{C_{\rm sb}}d_j^2(R_j^{\rm sb})^{-2M_{\rm sb}-6},
\]
which, since $d_j,R_j^{\rm sb}\ge1$, implies the deliberately weaker bound
\begin{equation}
\boxed{
\|\mathcal K_{j,\rm hi}\|_{S_2}^2
\ll L^{C_{\rm sb}}d_j^3(R_j^{\rm sb})^{-M_{\rm sb}}.
}
\label{eq:high-sideband-one-shell}
\end{equation}
Thus, since $d_j\asymp T/J_j$,
\begin{align*}
\sum_j\|\mathcal K_{j,\rm hi}\|_{S_2}^2
&\ll L^{C_{\rm sb}+\sigma M_{\rm sb}}
T^{3-M_{\rm sb}}
\sum_{J_j\le X_T}J_j^{M_{\rm sb}-3}\\
&\ll T L^{C_{\rm sb}+\sigma M_{\rm sb}
-A_0(M_{\rm sb}-2)}.
\end{align*}
By the prior choice \eqref{eq:acyclic-A0-choice}, the logarithmic exponent is
less than $-3$, so the right side is $o(TL)=o(N_T)$. No derivative order is
chosen at this stage. This piece has no diagonal jump because the local symbol
is $C^\infty$.

\emph{(viii) No local holomorphic/Perron complement.}
The final route never moves an exact HLP Dirichlet series to a local Perron
representation.  Hence there is no exact low-level holomorphic complement and
no levelwise Perron connector in the source ledger.  This class is identically
zero.

\emph{(ix) Finite-model scalar HLZ contour error.}
This class is controlled at the operator level by
Lemma~\ref{lem:HLZ-diagonal-operator-lift}.  The common-core congruence of
Proposition~\ref{prop:common-core-congruence} and the Fubini lift of
Lemma~\ref{lem:model-HLZ-Fubini-lift} place the packet polarization in the
analysis map before the scalar error is applied. Thus, with $R_T^{\rm com}$
denoting only the finite-model HLZ diagonalization error,
\[
\left\|\Gfr_T^{-1/2}R_T^{\rm com}\Gfr_T^{-1/2}\right\|_{S_2}^2
\ll N_T\varepsilon_{\rm diag}^2=o(N_T).
\]
There is no levelwise exact-HLP Perron contour error in the final route.

\emph{(x) Terminal logarithmic tail.}
The nonstationary part of the complementary physical band $W_T<u\le L$ is
exactly $R_T^{\rm tail}$ of Proposition~\ref{prop:terminal-log-tail}, and
already satisfies the required relative Hilbert--Schmidt bound after the
explicit one-leaking-leg lift. Its stationary and endpoint parts belong only
to $E_T^{\rm edge}$.

\emph{(xi) Exterior packet reserve and transition caps.}
The complete Fourier support of the actual smooth packet extends beyond the
central physical band by the two logarithmic reserves and the fixed-width
transition caps. Proposition~\ref{prop:exterior-packet-support} partitions all
terms with at least one such exterior leg as
\[
A_T^{\rm ext}=E_T^{\rm ext}+R_T^{\rm ext}.
\]
The stationary exterior part $E_T^{\rm ext}$ belongs only to
$E_T^{\rm edge}$, while
\[
\|\Gfr_T^{-1/2}R_T^{\rm ext}\Gfr_T^{-1/2}\|_{S_2}^2=o(N_T).
\]
Thus neither the flat reserve nor the smooth transition caps are omitted from
the source ledger.

Class (i) is absent and class (viii) vanishes. For (ii)--(v) and the smooth
height-partition corrections in (vi), the squared relative Hilbert--Schmidt
contribution is $L^{C_{\rm reg}+o(1)}$ per shell, so the total is
\[
\ll X_TL^{C_{\rm reg}+o(1)}
=T L^{C_{\rm reg}-A_0+o(1)}=o(N_T)
\]
by \eqref{eq:acyclic-A0-choice}. The nonstationary interior of (vi) is accounted
for by \eqref{eq:nonstat-relative-HS}; its endpoint part belongs instead to
$E_T^{\rm bdry}$ by Proposition~\ref{prop:boundary-endpoint-rank}. Class (vii)
has \eqref{eq:high-sideband-one-shell} and its global sum, class (ix) has
\eqref{eq:HLZ-diagonal-relative-HS}, and classes (x)--(xi) have
Propositions~\ref{prop:terminal-log-tail} and
\ref{prop:exterior-packet-support}. The primitive-controlled jump bounds in
(ii)--(v) are uniformly $o(1)$, so Lemma~\ref{lem:two-variable-primitive}
applies. Adding these classes together with (vi), (vii), (ix), (x), and (xi)
gives squared Hilbert--Schmidt norm $o(N_T)$, proving
\eqref{eq:regular-HS}.
\end{proof}

\begin{proposition}[Finite-$T$ decomposition of the contour form]
\label{prop:source-ledger}
For the normalized contour matrix $A_T$ of \eqref{eq:contour-form}, let $P_T$ be
Definition~\ref{def:PT}, let $E_T^{\rm edge}$ be the explicitly assembled edge
matrix of \eqref{eq:total-edge-rank}, and define the residual operator by the
exact algebraic difference
\begin{equation}
\boxed{R_T^{\rm tri}:=A_T-P_T-E_T^{\rm edge}.}
\label{eq:tri-remainder-definition}
\end{equation}
Then, identically,
\begin{equation}
\boxed{A_T=P_T+E_T^{\rm edge}+R_T^{\rm tri}.}
\label{eq:source-ledger}
\end{equation}
Moreover the sequential source splits in the proof below identify this
$R_T^{\rm tri}$ with the mutually exclusive regular classes of
Proposition~\ref{prop:regular-ledger}; hence, after quadratic-form compression
to $\Pret$ of
\eqref{eq:retained-space},
\begin{align}
P_T|_{\Pret}&\succeq c_0\,\Gfr_T|_{\Pret},\notag\\
\rank(E_T^{\rm edge}|_{\Pret})
&\le \rank E_T^{\rm edge}=o(N_T),\notag\\
\left\|
(\Gfr_T|_{\Pret})^{-1/2}(R_T^{\rm tri}|_{\Pret})
(\Gfr_T|_{\Pret})^{-1/2}
\right\|_{S_2}^2
&=o(N_T).
\end{align}
The point of the definition \eqref{eq:tri-remainder-definition} is that the
operator identity is not inferred from an informal ledger.  What remains to be
proved is the identification of this algebraic difference with the estimated
regular classes.  The ownership order is fixed as follows.  On the original
absolute-convergence right edge use the exact source split
\eqref{eq:direct-model-source-split}.  In the finite model, first perform the
residue-preserving local-height HLZ rescaling and the critical-$J$
smooth-to-sharp transfer.  The resulting sharp completed model is then split
exactly as $P_T+R_T^{\rm par}$ by
\eqref{eq:common-core-principal-congruence}.  Thus only the fully frozen part is
assigned to $P_T$; the parameter defect and the AFE smoothing discrepancy
belong only to regular class~(v).  The entire exact--model residual remains a
direct physical one-$\chi$ family; its low sidebands belong to
$E_T^{\rm edge}$ and its high sidebands to regular class~(vii).  The scalar self
difference $f-F_b$ belongs only to regular class~(ii), while the finite-model
scalar HLZ error belongs only to class~(ix).  No exact HLP local projector,
exact/model Laurent telescope, or levelwise Perron connector is used.
Independently, the complete smooth packet support is first split into the
central physical band and the two exterior reserve/cap regions of
Proposition~\ref{prop:exterior-packet-support}.  The central physical band is
then split into the stable core and its terminal complement.  The stationary
exterior contribution belongs to $E_T^{\rm edge}$ and its nonstationary
complement only to regular class~(xi). By
Proposition~\ref{prop:trivial-outer-twist}, there is no additional outer
arithmetic annulus in the invariant source.
\end{proposition}

\begin{proof}
The identity \eqref{eq:source-ledger} itself is immediate from the definition
\eqref{eq:tri-remainder-definition}.  We verify that this algebraic remainder
is exactly the sum estimated in Proposition~\ref{prop:regular-ledger}.

Start with the exact same-form identity \eqref{eq:same-form-identity} and the
matrix-level right/left decomposition \eqref{eq:J-I-source-split}.  On the
right edge, \eqref{eq:exact-fprime-separation} is an algebraic identity: the
principal factor $\mathcal C\widehat H_{\rm ex}$ and the $f'$ source are
separated without approximation.  Horizontal and nonstationary pieces are
then separated by the literal contour and stationary/nonstationary partitions,
so their sum with the stationary right-edge source is still the original
matrix.

On each retained stable height block use the exact identity
\eqref{eq:direct-model-source-split}.  For the finite model term, the self face
is treated by the critical-$J$ identity and the logarithmic-derivative faces by
the scalar three-$Z$ master.  Lemma~\ref{lem:HLZ-local-height-rescaling} first
puts the three-$Z$ diagonal on the local physical scale, and
Lemma~\ref{lem:J-AFE-sharp-transfer} splits the critical-$J$ face into its sharp
full fibre plus its smoothing discrepancy.  Proposition~\ref{prop:common-core-congruence}
then splits the sharp completed model as $P_T+R_T^{\rm par}$; the former is the
fully frozen positive term, while $R_T^{\rm par}$ and the AFE discrepancy are
assigned exactly once to regular class~(v).  Lemma~\ref{lem:model-HLZ-Fubini-lift}
keeps the scalar HLZ contour error of the three-$Z$ faces in regular class~(ix).
The scalar difference $f-F_b$ is assigned exactly once to class~(ii).

The remaining source is precisely $\Delta H_{\rm dir}$.  By
Proposition~\ref{prop:direct-exact-model-residual} its complete HLP hierarchy
and the two $P$-terms retain the original direct physical carrier.  Apply the
literal disjoint sideband indicators $|r|\le R_j^{\rm sb}$ and
$|r|>R_j^{\rm sb}$ to this whole norm-convergent family.  The low sidebands
belong exactly once to $E_T^{\rm edge}$, while the high sidebands belong exactly
once to regular class~(vii).  No exact HLP term is also represented inside
$P_T$.  The endpoint integrations by parts likewise split into their finite jet
terms and an interior remainder, assigning the former exactly once to
$E_T^{\rm bdry}$.

Independently split the complete smooth packet support into the exterior
reserve/cap region and the central physical band.  Proposition~\ref{prop:exterior-packet-support}
partitions the former into its stationary edge part $E_T^{\rm ext}$ and its
nonstationary regular part $R_T^{\rm ext}$, including all cross terms with one
exterior leg.  On the central physical band split the logarithmic coordinate at
$W_T$. The core piece is the one used in $P_T$ and in the preceding stable-band
decomposition; Proposition~\ref{prop:terminal-log-tail} partitions the
complementary band into its stationary edge, boundary, and nonstationary
regular pieces, including the core--tail cross terms stated there.
Proposition~\ref{prop:trivial-outer-twist} sets the auxiliary outer indices to
$1$, so no additional arithmetic annulus is left over.

Every step above is either an exact algebraic identity, a contour identity with
its connector term retained, or a disjoint index/support partition.  Hence the
pieces assigned to $P_T$ and $E_T^{\rm edge}$ leave precisely the eleven
regular classes enumerated in Proposition~\ref{prop:regular-ledger}.  By
\eqref{eq:tri-remainder-definition} their sum is $R_T^{\rm tri}$.  The three
estimates in the proposition are therefore
\eqref{eq:principal-floor}, \eqref{eq:total-edge-rank}, and
\eqref{eq:regular-HS}, respectively.
\end{proof}
